\subsection*{逐段走查核函数}

阶段索引为 \co@iv_k % STAGES@。使用两个槽位时，分块 0 和 2 使用槽位 0，分块 1 和 3 使用槽位 1。它们是不同的逻辑分块，在不同时刻共用物理存储；事件对属于槽位，不属于某个永久分块。

生产者在写入任一操作数前等待 \co@empty[stage]@。本例的两次复制均为同步操作，因此在它们之后触发 \co@full[stage]@，就表示两个操作数都已就绪。若替换为异步搬运，则必须先完成这些搬运，再发出同一个触发。

消费者先触发所有 \co@empty@ 事件。事件初始均未置位，因此这一步授予生产者首次填充各槽位的许可。此时尚未装入输入，所以不会触发任何 \co@full@ 事件。

消费者等待 \co@full[stage]@，装入操作数片段并累加当前切片，随后触发 \co@empty[stage]@。这个顺序不能颠倒：发出矩阵操作，并不足以证明它已经停止读取共享操作数。

完成点保护缓冲区复用。若 MMA 仍在读取时生产者就覆盖槽位，即使分块形状完全正确，消费者也可能混用两个 K 切片的操作数。仅检查 \co@full@ 事件无法防止这一错误，归还 \co@empty@ 事件同样必要。
